
\documentclass[11pt,a4paper]{article}

% Template ACL chinh thuc
\usepackage[final]{acl}

% Ho tro tieng Viet voi XeLaTeX
\usepackage{fontspec}
\setmainfont{Noto Serif}
\usepackage{microtype}

\title{
Dịch máy Anh--Việt sử dụng Seq2Seq LSTM với Attention
và Transformer: Đề xuất nghiên cứu
}

% THAY HO TEN VA MSSV TRUOC KHI NOP
\author{
Tú (MSSV: điền mã) \quad
Trung (MSSV: điền mã) \quad
Đức (MSSV: điền mã) \\
Nhóm: điền mã nhóm
}

\begin{document}
\maketitle

\begin{abstract}
Đề tài nghiên cứu bài toán dịch máy thần kinh từ
tiếng Anh sang tiếng Việt bằng hai kiến trúc:
Seq2Seq LSTM kết hợp cơ chế Attention và
Transformer Encoder--Decoder. Nhóm sử dụng bộ
dữ liệu song ngữ PhoMT, xây dựng quy trình tiền
xử lý và huấn luyện thống nhất, sau đó so sánh
chất lượng bản dịch bằng corpus BLEU. Nghiên cứu
tập trung vào việc tự cài đặt các thành phần
cốt lõi, phân tích lỗi dịch và trực quan hóa
attention. Do hạn chế về tài nguyên tính toán,
nhóm dự kiến tiến hành thí nghiệm trên tập con
của PhoMT bằng GPU miễn phí.
\end{abstract}

\section{Giới thiệu}

Dịch máy thần kinh (Neural Machine Translation,
NMT) sử dụng mạng học sâu để ánh xạ một chuỗi
ngôn ngữ nguồn thành chuỗi ngôn ngữ đích.
Seq2Seq là một hướng tiếp cận nền tảng, trong
khi Attention giúp mô hình khai thác thông tin
từ các vị trí khác nhau của câu nguồn
\citep{bahdanau2015nmt}. Transformer thay thế
các bước tính toán hồi quy bằng cơ chế
self-attention \citep{vaswani2017attention}.

Mục tiêu của nhóm là tự cài đặt và so sánh
hai kiến trúc trên cùng bài toán dịch
\textbf{Anh sang Việt}, qua đó nghiên cứu
sự khác biệt về chất lượng dịch, tốc độ
huấn luyện và các dạng lỗi thường gặp.

\section{Dữ liệu}

Nhóm sử dụng PhoMT, bộ dữ liệu song ngữ
Việt--Anh quy mô lớn được giới thiệu bởi
\citet{doan2021phomt}. Bộ dữ liệu có
3.015.869 cặp câu, gồm 2.977.999 cặp train,
18.719 cặp validation và 19.151 cặp test.

Nhóm giữ nguyên cách chia dữ liệu được
công bố, chỉ sử dụng tập train để xây dựng
tokenizer và từ vựng. Các bước chuẩn bị
dữ liệu gồm kiểm tra cặp câu, lọc câu không
hợp lệ, chuẩn hóa văn bản và chuyển câu
thành chuỗi token. Tập con huấn luyện
sẽ được chọn từ train với seed cố định
để phù hợp tài nguyên GPU miễn phí.

\section{Phương pháp}

\paragraph{Seq2Seq LSTM với Attention.}
Mô hình gồm encoder LSTM biểu diễn câu nguồn
và decoder LSTM sinh từng token tiếng Việt.
Cơ chế Attention tính trọng số trên các
hidden state của encoder để hình thành
context vector tại mỗi bước giải mã.
Nhóm tự xây dựng các thành phần encoder,
decoder và Attention.

\paragraph{Transformer.}
Nhóm tự xây dựng Transformer nhỏ theo
\citet{vaswani2017attention}, gồm positional
encoding, multi-head attention, feed-forward
network, residual connection và layer
normalization. Decoder sử dụng causal mask
để không nhìn thấy các token tương lai.

\paragraph{Đối chứng.}
Hai mô hình được huấn luyện và đánh giá trên
cùng tập dữ liệu và cùng quy trình tiền xử lý.
Các cấu hình và ngân sách tính toán sẽ
được ghi nhận để việc so sánh minh bạch.

\section{Thiết kế thực nghiệm}

Nhóm dự kiến sử dụng PyTorch và GPU miễn phí
trên Kaggle hoặc Google Colab. Quá trình
huấn luyện sẽ theo dõi training loss,
validation loss, thời gian mỗi epoch và
lưu checkpoint của mô hình tốt nhất trên
tập validation.

Chất lượng dịch được đánh giá bằng
\textbf{corpus BLEU} \citep{papineni2002bleu}.
Nhóm sử dụng cùng tập test, cùng phương pháp
giải mã và cấu hình đánh giá được công bố
rõ ràng. Ngoài BLEU, nhóm phân tích các
trường hợp dịch thiếu từ, sai nghĩa và
câu dài; đồng thời trực quan hóa
attention map.

\section{Kế hoạch và phân công}

\textbf{Tuần 1:} Khảo sát bài báo, kiểm tra
PhoMT, chuẩn bị môi trường và Proposal.

\textbf{Tuần 2:} Hoàn thiện pipeline dữ liệu,
tokenizer, baseline và huấn luyện thử.

\textbf{Tuần 3:} Hoàn thiện hai mô hình,
checkpoint và báo cáo tiến độ.

\textbf{Tuần 4:} Huấn luyện, so sánh BLEU và
thực hiện ablation.

\textbf{Tuần 5:} Phân tích lỗi, trực quan hóa
attention và soạn báo cáo.

\textbf{Tuần 6:} Hoàn thiện báo cáo ACL,
slide và demo; chuẩn bị bảo vệ cá nhân.

\textbf{Tú:} Dữ liệu, tokenizer, DataLoader,
BLEU và tổng hợp Proposal.

\textbf{Trung:} Seq2Seq Attention, giải mã,
demo và báo cáo tiến độ.

\textbf{Đức:} Transformer, huấn luyện,
checkpoint, logging và tổng hợp báo cáo cuối.

Các thành viên cùng tham gia kiểm thử,
đánh giá kết quả, review Pull Request và
chuẩn bị trả lời vấn đáp.

\section{Kết luận}

Đề tài hướng đến một hệ thống dịch máy
Anh--Việt có mã nguồn tự triển khai,
quy trình thực nghiệm tái lập được và
phân tích có căn cứ về ưu, nhược điểm
của Seq2Seq Attention và Transformer.

\bibliography{references}

\end{document}
